\documentclass[a4paper]{article}
\usepackage[pdftex]{graphicx}
\usepackage[utf8]{inputenc}
\usepackage{enumerate}
\usepackage{amssymb}
\usepackage{icomma}
\usepackage{tikz}
\usepackage{href-ul}
\hypersetup{
	colorlinks=true,
	linkcolor=blue,
	urlcolor=blue}
\usepackage{geometry}
\geometry{a4paper, top=15mm, left=15mm, right=15mm, bottom=15mm,
	headsep=10mm, footskip=12mm}

% Start the document
\begin{document}
\begin{tikzpicture}(10,3)
	\draw[ultra thick](2,0) --(10,0) -- (10,1.3) --(2,1.3) -- (2,0);
	\draw[fill=black](2,0)-- (0,.6) -- (2,.6) -- (2,0);
	\node at (6,.6) {\large Lösungsblatt von \href{https://www.okuyakl.de}{www.okuyakl.de}};
\end{tikzpicture}
\vspace{0.5 cm}


\renewcommand{\arraystretch}{3}
\begin{tabular}{|p{40 pt}|p{80 pt}|p{100 pt}|p{100 pt}|p{80 pt}|}
\hline
Gerade & Steigung $m$ & Y-Achsenabschnitt $t$  & Funktionsterm $y=$ & Nullstelle $x=$\\
\hline
$g_1:$ & $-4$ & 5 & $y=-4x+5$ & 1,25 \\
\hline
$g_2:$ & $-{1\over 3} $& 3 & $y=-{1\over 3}x+3$  & 9 \\
\hline
$g_3:$ & ${1\over 9}$ & 2 & $y={1\over 9}x+2$ & $-18$\\
\hline
$g_4:$ & 1 & 0 & $y=x$ & 0 \\
\hline
$g_5:$ & $-0,25$ & $-2$ & $y=-0,25x-2$ & $-8$\\
\hline
$g_6:$ & 2 & $-3$ & $y=2x-3 $ & $1,5$\\
\hline
$g_7:$ &$ -{2\over 3}$& $-3$ & $y=-{2\over 3}x-3$ & $-4,5$ \\
\hline
\end{tabular}

\begin{center}
	\includegraphics[width=7 cm]{../../viecher/endcomic.pdf}
	
	Hier geht es zurück zum \href{https://www.okuyakl.de/math/m8gegeL017/aa017.pdf}{Aufgabenblatt}
\end{center}

\end{document}